\documentclass[12pt]{article}
\usepackage[T1]{fontenc}
\usepackage[utf8]{inputenc}
\usepackage{lmodern}
\usepackage{textcomp}
\usepackage{enumitem}
\usepackage{geometry}
\geometry{margin=1in}
\begin{document}
\begin{center}
Answer Key - Biology - Graduate Level Assessment 24 \
\small (Answers are ordered exactly as in the question paper.)
\end{center}

\section*{Multiple Choice}
\begin{enumerate}[label=\arabic*.]
\item \textbf{Answer:} A
\\ \textit{Options:} A) the tRNA transcript from the original DNA; B) a single strand of the original DNA segment after a substitution mutation; C) the final processed mRNA made from the original DNA; D) the primary RNA transcript (after splicing) from the original DNA; E) a single strand of the original DNA segment after a duplication mutation
\\ \textit{Note:} The tRNA transcript from the original DNA contains the fewest number of nucleotides because tRNA is a smaller molecule that is produced from a single gene segment and undergoes processing, whereas mRNA and rRNA are generally longer and contain more nucleotides due to their coding and structural roles.
\item \textbf{Answer:} A
\\ \textit{Options:} A) Plants, animals, and fungi all have both chloroplasts and mitochondria.; B) Plants and fungi have chloroplasts but no mitochondria; animals have only mitochondria.; C) Animals have both chloroplasts and mitochondria; plants and fungi have only chloroplasts.; D) Plants and animals have chloroplasts but no mitochondria; fungi have only mitochondria.; E) Plants and animals have mitochondria but no chloroplasts; fungi have only chloroplasts.
\\ \textit{Note:} The correct answer is misleading because while plants possess both chloroplasts for photosynthesis and mitochondria for cellular respiration, animals and fungi only have mitochondria, highlighting the distinction in organelle presence across different kingdoms of life.
\item \textbf{Answer:} A
\\ \textit{Options:} A) hydrogen; B) methane; C) nitrogen; D) sulfur dioxide; E) neon
\\ \textit{Note:} The correct answer is hydrogen because while primitive Earth's atmosphere is thought to have contained gases like methane, ammonia, and water vapor, free molecular hydrogen was likely scarce due to its lightness and tendency to escape the planet's gravity.
\item \textbf{Answer:} A
\\ \textit{Options:} A) Enzymatic reactions with oxygen; B) Lactose crystallization; C) Loss of refrigeration power over time; D) mesophiles; E) Natural separation of milk components without bacterial influence
\\ \textit{Note:} Milk spoils in the refrigerator primarily due to enzymatic reactions with oxygen, which catalyze the breakdown of components in the milk, leading to off-flavors and spoilage despite the cold temperatures slowing down microbial growth.
\item \textbf{Answer:} D
\\ \textit{Options:} A) Cue, Desire, Fulfillment; B) Arousal, Behavior, Consequence; C) Stimulus, Habituation, Extinction; D) Appetitive behavior, Consummatory act, Quiescence; E) Perception, Decision, Movement
\\ \textit{Note:} The correct answer reflects the three typical phases of a behavioral act, where appetitive behavior represents the seeking or motivation for a behavior, the consummatory act is the culmination of that behavior, and quiescence indicates a state of rest or recovery following the behavior.
\end{enumerate}

\section*{True / False}
\begin{enumerate}[label=\arabic*.]
\setcounter{enumi}{5}
\item \textbf{Answer:} True
\\ \textit{Options:} A) True; B) False
\\ \textit{Note:} Bedtime NPH insulin added to maximal therapy with sulfonylurea and metformin is an effective, simple, well-tolerated approach for patients with uncontrolled type 2 diabetes.
\item \textbf{Answer:} True
\\ \textit{Options:} A) True; B) False
\\ \textit{Note:} Human pWAT has chemotactic properties through the secretion of different chemokines, and we propose that pWAT might contribute to the progression of obesity-associated atherosclerosis.
\item \textbf{Answer:} False
\\ \textit{Options:} A) True; B) False
\\ \textit{Note:} There is little evidence that the 2008 public antibiotic campaigns were effective. The use and visibility of future campaign materials needs auditing. A carefully planned approach that targets the public in GP waiting rooms and through clinicians in consultations may be a more effective way of improving prudent antibiotic use.
\item \textbf{Answer:} True
\\ \textit{Options:} A) True; B) False
\\ \textit{Note:} The real-time PCR approach revealed promising results in pollen identification and quantification, even when analyzing pollen mixes. Future perspectives could concern the development of multiplex real-time PCR for the simultaneous detection of different taxa in the same reaction tube and the application of high-throughput molecular methods.
\item \textbf{Answer:} False
\\ \textit{Options:} A) True; B) False
\\ \textit{Note:} The new storage can affords more stable temperature levels when compared to the formerly used can. Since temperature is stable during conservation, continuous monitoring in everyday practice does not seem warranted.
\end{enumerate}

\section*{Long Form Response}
\begin{enumerate}[label=\arabic*.]
\setcounter{enumi}{10}
\item \textbf{Answer:} Certain animals, such as birds and reptiles, employ various physiological and behavioral mechanisms to regulate their body temperature despite environmental fluctuations. Birds, being endothermic, maintain a stable internal temperature through metabolic heat production, enhanced by their high metabolic rates and insulating feathers, which minimize heat loss. In contrast, reptiles, which are ectothermic, rely on behavioral adaptations such as basking in sunlight to increase their body temperature or seeking shade to cool down, effectively using their environment to thermoregulate. Additionally, both groups exhibit physiological responses; for example, birds can alter blood flow to their extremities to conserve or dissipate heat, while reptiles may change their body posture to optimize heat absorption. These adaptations illustrate the diverse strategies employed by different taxa to achieve thermoregulation in varying ecological contexts.
\\ \textit{Note:} The answer correctly identifies the divergent thermoregulatory strategies of birds and reptiles, emphasizing the role of metabolic processes in endothermic birds and behavioral adaptations in ectothermic reptiles to maintain body temperature amidst environmental changes.
\item \textbf{Answer:} Stabilizing selection is exemplified in Swiss starlings through their reproductive patterns, particularly regarding clutch size, which refers to the number of eggs laid in a single nesting attempt. Research indicates that intermediate clutch sizes yield the highest reproductive success, as excessively small clutches fail to maximize the number of offspring, while overly large clutches can lead to inadequate care and lower survival rates. This balancing act reflects the selective pressure favoring starlings that produce a moderate number of eggs, thus enhancing their overall fitness. Consequently, stabilizing selection plays a crucial role in maintaining the optimal clutch size within the population, ensuring that the traits associated with reproductive success are preserved over generations.
\\ \textit{Note:} correct because it illustrates how stabilizing selection favors intermediate clutch sizes in Swiss starlings, as these sizes optimize reproductive success by balancing the number of offspring with the ability to care for them, thereby enhancing fitness.
\end{enumerate}

\vfill
\noindent\textit{End of Answer Key}
\end{document}